\chapter{2020年全国II卷（文）}

\section{单选题}

\begin{problem}
已知集合 \(A=\{x\mid \lvert x\rvert<3,\ x\in \Z\}\)，\(B=\{x\mid \lvert x\rvert>1,\ x\in \Z\}\)，则 \(A\cap B=\)（\quad）

\choices
    {\(\varnothing\)}
    {\(\{-3,-2,2,3\}\)}
    {\(\{-2,0,2\}\)}
    {\(\{-2,2\}\)}
\end{problem}

\begin{answer}
D
\end{answer}

\begin{solution}
由 \(\lvert x\rvert<3\) 且 \(x\in\Z\)，得 \(A=\{-2,-1,0,1,2\}\)；由 \(\lvert x\rvert>1\) 且 \(x\in\Z\)，得 \(B\) 中满足条件的整数的绝对值至少为 \(2\). 因此 \(A\cap B=\{-2,2\}\)，故选 D.
\end{solution}

\begin{problem}
\((1-\i)^4=\)（\quad）

\choices
    {\(-4\)}
    {\(4\)}
    {\(-4\i\)}
    {\(4\i\)}
\end{problem}

\begin{answer}
A
\end{answer}

\begin{solution}
先计算平方得 \((1-\i)^2=1-2\i+\i^2=-2\i\)，所以 \((1-\i)^4=(-2\i)^2=4\i^2=-4\)，故选 A.
\end{solution}

\begin{problem}
如图，将钢琴上的 \(12\) 个键依次记为 \(a_1\)，\(a_2\)，\(\cdots\)，\(a_{12}\). 设 \(1\le i<j<k\le 12\). 若 \(k-j=3\) 且 \(j-i=4\)，则 \(a_i\)，\(a_j\)，\(a_k\) 为原位大三和弦；若 \(k-j=4\) 且 \(j-i=3\)，则称 \(a_i\)，\(a_j\)，\(a_k\) 为原位小三和弦. 用这 \(12\) 个键可以构成的原位大三和弦与原位小三和弦的个数之和为（\quad）

\FigureLayoutDeclare{img/2020/national_paper_2_liberal/q3_fig1.png}{14.0cm}{17bp 19bp 31bp 22bp}
\bitmapfigure[max width=6.0cm]{img/2020/national_paper_2_liberal/q3_fig1.png}

\choices
    {\(5\)}
    {\(8\)}
    {\(10\)}
    {\(15\)}
\end{problem}

\begin{answer}
C
\end{answer}

\begin{solution}
对于原位大三和弦，由 \(j=i+4\)，\(k=i+7\)，且 \(k\leq12\)，得 \(1\leq i\leq5\)，所以共有 \(5\) 个. 对于原位小三和弦，由 \(j=i+3\)，\(k=i+7\)，同样得 \(1\leq i\leq5\)，所以共有 \(5\) 个. 因此个数之和为 \(5+5=10\)，故选 C.
\end{solution}

\begin{problem}
在新冠肺炎疫情防控期间，某超市开通网上销售业务，每天能完成 \(1200\) 份订单的配货，由于订单量大幅增加，导致订单积压. 为解决困难，许多志愿者踊跃报名参加配货工作. 已知该超市某日积压 \(500\) 份订单未配货，预计第二天的新订单超过 \(1600\) 份的概率为 \(0.05\). 志愿者每人每天能完成 \(50\) 份订单的配货，为使第二天完成积压订单及当日订单的配货的概率不小于 \(0.95\)，则至少需要志愿者（\quad）

\choices
    {\(10\text{ 名}\)}
    {\(18\text{ 名}\)}
    {\(24\text{ 名}\)}
    {\(32\text{ 名}\)}
\end{problem}

\begin{answer}
B
\end{answer}

\begin{solution}
设第二天的新订单数为 \(X\)，志愿者人数为 \(n\). 第二天需要完成的订单总数为 \(500+X\)，而可完成的订单数为 \(1200+50n\). 当 \(X\leq1600\) 时，若要保证完成所有订单，需有 \(1200+50n\geq500+1600=2100\). 由于 \(P(X>1600)=0.05\)，故 \(P(X\leq1600)=0.95\)，上述容量条件正好保证完成概率不小于 \(0.95\). 于是 \(1200+50n\geq2100\)，解得 \(n\geq18\). 因而至少需要 \(18\) 名志愿者，故选 B.
\end{solution}

\begin{problem}
已知单位向量 \(\bs a\)，\(\bs b\) 的夹角为 \(60^\circ\)，则在下列向量中，与 \(\bs b\) 垂直的是（\quad）

\choices
    {\(\bs a+2\bs b\)}
    {\(2\bs a+\bs b\)}
    {\(\bs a-2\bs b\)}
    {\(2\bs a-\bs b\)}
\end{problem}

\begin{answer}
D
\end{answer}

\begin{solution}
因为 \(\lvert\bs a\rvert=\lvert\bs b\rvert=1\)，且夹角为 \(60^\circ\)，所以 \(\bs a\cdot\bs b=\cos60^\circ=\dfrac12\)，并且 \(\bs b\cdot\bs b=1\). 分别计算各选项与 \(\bs b\) 的数量积：
\(\left(\bs a+2\bs b\right)\cdot\bs b=\dfrac12+2\)，\(\left(2\bs a+\bs b\right)\cdot\bs b=1+1\)，\(\left(\bs a-2\bs b\right)\cdot\bs b=\dfrac12-2\)，\(\left(2\bs a-\bs b\right)\cdot\bs b=1-1=0\). 因此只有 \(2\bs a-\bs b\) 与 \(\bs b\) 垂直，故选 D.
\end{solution}

\begin{problem}
记 \(S_n\) 为等比数列 \(\{a_n\}\) 的前 \(n\) 项和. 若 \(a_5-a_3=12\)，\(a_6-a_4=24\)，则 \(\dfrac{S_n}{a_n}=\)（\quad）

\choices
    {\(2^n-1\)}
    {\(2-2^{1-n}\)}
    {\(2-2^{n-1}\)}
    {\(2^{1-n}-1\)}
\end{problem}

\begin{answer}
B
\end{answer}

\begin{solution}
设等比数列的公比为 \(q\). 由 \(a_5-a_3=12\ne0\)，可知相关因子不为零，并且
\(\dfrac{a_6-a_4}{a_5-a_3}=\dfrac{a_1q^5-a_1q^3}{a_1q^4-a_1q^2}=q\). 因此 \(q=\dfrac{24}{12}=2\). 于是
\(S_n=a_1\dfrac{2^n-1}{2-1}=a_1(2^n-1)\)，\(a_n=a_1\cdot2^{n-1}\)，从而 \(\dfrac{S_n}{a_n}=\dfrac{2^n-1}{2^{n-1}}=2-2^{1-n}\)，故选 B.
\end{solution}

\begin{problem}
执行下面的程序框图，若输入的 \(k=0\)，\(a=0\)，则输出的 \(k\) 为（\quad）

\bitmapinlinefigure[5.0cm][width=5.0cm]{img/2020/national_paper_2_liberal/q7_fig1.png}

\choices
    {\(2\)}
    {\(3\)}
    {\(4\)}
    {\(5\)}
\end{problem}

\begin{answer}
C
\end{answer}

\begin{solution}
程序开始时 \(k=0\)，\(a=0\). 第一次循环后 \(a=2\times0+1=1\)，\(k=1\)；第二次循环后 \(a=3\)，\(k=2\)；第三次循环后 \(a=7\)，\(k=3\)；第四次循环后 \(a=15\)，\(k=4\). 此时 \(a>10\)，输出 \(k=4\)，故选 C.
\end{solution}

\begin{problem}
若过点 \((2,1)\) 的圆与两坐标轴都相切，则圆心到直线 \(2x-y-3=0\) 的距离为（\quad）

\choices
    {\(\dfrac{\sqrt5}{5}\)}
    {\(\dfrac{2\sqrt5}{5}\)}
    {\(\dfrac{3\sqrt5}{5}\)}
    {\(\dfrac{4\sqrt5}{5}\)}
\end{problem}

\begin{answer}
B
\end{answer}

\begin{solution}
设圆心为 \((h,k)\)，半径为 \(r>0\). 圆与两坐标轴都相切，所以 \(\lvert h\rvert=\lvert k\rvert=r\). 又因为圆过点 \((2,1)\)，有
\((2-h)^2+(1-k)^2=r^2\)，即 \(r^2-4h-2k+5=0\). 分别令 \(h=\varepsilon r\)，\(k=\delta r\)，其中 \(\varepsilon,\delta\in\{1,-1\}\)，得到 \(r^2-(4\varepsilon+2\delta)r+5=0\). 当 \((\varepsilon,\delta)=(1,1)\) 时，方程为 \(r^2-6r+5=0\)，得 \(r=1\) 或 \(r=5\)；当 \((\varepsilon,\delta)=(1,-1)\) 或 \((-1,1)\) 时，方程分别为 \(r^2-2r+5=0\) 或 \(r^2+2r+5=0\)，二者的判别式均为 \(-16<0\)，均无实根；当 \((\varepsilon,\delta)=(-1,-1)\) 时，方程为 \(r^2+6r+5=0\)，得 \(r=-1\) 或 \(r=-5\)，均不满足 \(r>0\). 因此圆心只能是 \((1,1)\) 或 \((5,5)\). 对这两个圆心，圆心到直线 \(2x-y-3=0\) 的距离都为 \(\dfrac{\lvert2h-k-3\rvert}{\sqrt{2^2+(-1)^2}}=\dfrac{2}{\sqrt5}=\dfrac{2\sqrt5}{5}\)，故选 B.
\end{solution}

\begin{problem}
设 \(O\) 为坐标原点，直线 \(x=a\) 与双曲线 \(C:\dfrac{x^2}{a^2}-\dfrac{y^2}{b^2}=1\)（\(a>0,\ b>0\)）的两条渐近线分别交于 \(D\)，\(E\) 两点. 若 \(\triangle ODE\) 的面积为 \(8\)，则 \(C\) 的焦距的最小值为（\quad）

\choices
    {\(4\)}
    {\(8\)}
    {\(16\)}
    {\(32\)}
\end{problem}

\begin{answer}
B
\end{answer}

\begin{solution}
双曲线的两条渐近线为 \(y=\dfrac ba x\) 和 \(y=-\dfrac ba x\). 它们与直线 \(x=a\) 的交点为 \(D(a,b)\)，\(E(a,-b)\)，所以 \(\triangle ODE\) 的面积为 \(\dfrac12\cdot2b\cdot a=ab=8\). 双曲线的焦半距为 \(c=\sqrt{a^2+b^2}\)，焦距为 \(2c\). 由基本不等式 \(a^2+b^2\geq2ab=16\)，得 \(2c=2\sqrt{a^2+b^2}\geq2\sqrt{16}=8\). 当 \(a=b=2\sqrt2\) 时取等号，因此焦距的最小值为 \(8\)，故选 B.
\end{solution}

\begin{problem}
设函数 \(f(x)=x^3-\dfrac{1}{x^3}\)，则 \(f(x)\)（\quad）

\choices
    {是奇函数，且在 \((0,+\infty)\) 单调递增}
    {是奇函数，且在 \((0,+\infty)\) 单调递减}
    {是偶函数，且在 \((0,+\infty)\) 单调递增}
    {是偶函数，且在 \((0,+\infty)\) 单调递减}
\end{problem}

\begin{answer}
A
\end{answer}

\begin{solution}
函数定义域为 \(\R\smallsetminus\{0\}\). 对任意 \(x\ne0\)，有 \(f(-x)=(-x)^3-\dfrac{1}{(-x)^3}=-x^3+\dfrac1{x^3}=-f(x)\)，所以 \(f(x)\) 是奇函数. 又在 \((0,+\infty)\) 上，\(f'(x)=3x^2+\dfrac{3}{x^4}>0\)，故 \(f(x)\) 在该区间上单调递增，故选 A.
\end{solution}

\begin{problem}
已知 \(\triangle ABC\) 是面积为 \(\dfrac{9\sqrt3}{4}\) 的等边三角形，且其顶点都在球 \(O\) 的球面上. 若球 \(O\) 的表面积为 \(16\pi\)，则 \(O\) 到平面 \(ABC\) 的距离为（\quad）

\choices
    {\(\sqrt3\)}
    {\(\dfrac{3}{2}\)}
    {\(1\)}
    {\(\dfrac{\sqrt3}{2}\)}
\end{problem}

\begin{answer}
C
\end{answer}

\begin{solution}
设等边三角形边长为 \(s\). 由面积公式 \(\dfrac{\sqrt3}{4}s^2=\dfrac{9\sqrt3}{4}\)，得 \(s=3\). 等边三角形的外接圆半径为 \(R_{\mathrm{tri}}=\dfrac{s}{\sqrt3}=\sqrt3\). 设球 \(O\) 的半径为 \(R\)，由 \(4\pi R^2=16\pi\) 得 \(R=2\). 球心到平面 \(ABC\) 的垂足是三角形外心，故由直角三角形关系得 \(d^2+R_{\mathrm{tri}}^2=R^2\)，即 \(d^2+3=4\)，所以 \(d=1\)，故选 C.
\end{solution}

\begin{problem}
若 \(2^x-2^y<3^{-x}-3^{-y}\)，则（\quad）

\choices
    {\(\ln(y-x+1)>0\)}
    {\(\ln(y-x+1)<0\)}
    {\(\ln\lvert x-y\rvert>0\)}
    {\(\ln\lvert x-y\rvert<0\)}
\end{problem}

\begin{answer}
A
\end{answer}

\begin{solution}
令 \(t=y-x\). 原不等式化为 \(2^x(1-2^t)<3^{-x}(1-3^{-t})\). 当 \(t>0\) 时，左边为负数，右边为正数，不等式成立；当 \(t=0\) 时两边相等，不等式不成立；当 \(t<0\) 时，左边为正数，右边为负数，不等式不成立. 因而必有 \(y-x>0\)，所以 \(y-x+1>1\)，从而 \(\ln(y-x+1)>0\)，故选 A.
\end{solution}

\section{填空题}

\begin{problem}
若 \(\sin x=-\dfrac{2}{3}\)，则 \(\cos 2x=\fillinblank{}\).
\end{problem}

\begin{answer}
\(\dfrac{1}{9}\)
\end{answer}

\begin{solution}
由二倍角公式 \(\cos2x=1-2\sin^2x\)，代入 \(\sin x=-\dfrac23\)，得 \(\cos2x=1-2\times\dfrac49=\dfrac19\). 故答案为 \(\dfrac19\).
\end{solution}

\begin{problem}
记 \(S_n\) 为等差数列 \(\{a_n\}\) 的前 \(n\) 项和. 若 \(a_1=-2\)，\(a_2+a_6=2\)，则 \(S_{10}=\fillinblank{}\).
\end{problem}

\begin{answer}
\(25\)
\end{answer}

\begin{solution}
设等差数列的公差为 \(d\). 由 \(a_1=-2\)，得 \(a_2=-2+d\)，\(a_6=-2+5d\). 由 \(a_2+a_6=2\)，得 \(-4+6d=2\)，所以 \(d=1\). 因此 \(S_{10}=\dfrac{10}{2}\left(2a_1+9d\right)=5(-4+9)=25\). 故答案为 \(25\).
\end{solution}

\begin{problem}
若 \(x\)，\(y\) 满足约束条件
\(\begin{cases}
x+y\ge -1,\\
x-y\ge -1,\\
2x-y\le 1
\end{cases}\)，则 \(z=x+2y\) 的最大值是\fillinblank{}.
\end{problem}

\begin{answer}
\(8\)
\end{answer}

\begin{solution}
将约束条件改写为 \(y\geq-1-x\)，\(y\leq x+1\)，\(y\geq2x-1\). 可行域的边界交点分别由
\(y=-1-x\) 与 \(y=x+1\)，\(y=-1-x\) 与 \(y=2x-1\)，\(y=x+1\) 与 \(y=2x-1\) 求得，依次为 \((-1,0)\)，\((0,-1)\)，\((2,3)\). 目标函数 \(z=x+2y\) 在有界可行域的顶点处取得最大值，三个顶点处的值分别为 \(-1\)，\(-2\)，\(8\)，故最大值为 \(8\). 故答案为 \(8\).
\end{solution}

\begin{problem}
设有下列四个命题：

\begin{circlelist}
\item \(p_1\)：两两相交且不过同一点的三条直线必在同一平面内；
\item \(p_2\)：过空间中任意三点有且仅有一个平面；
\item \(p_3\)：若空间两条直线不相交，则这两条直线平行；
\item \(p_4\)：若直线 \(l\subset \alpha\)，直线 \(m\perp \alpha\)，则 \(m\perp l\).
\end{circlelist}

则下述命题中所有真命题的序号是\fillinblank{}

\begin{circlelist}
    \item \(p_1\land p_4\)
    \item \(p_1\land p_2\)
    \item \(\neg p_2\vee p_3\)
    \item \(\neg p_3\vee \neg p_4\)
\end{circlelist}
\end{problem}

\begin{answer}
\(\circled{1}\)\(\circled{3}\)\(\circled{4}\)
\end{answer}

\begin{solution}
命题 \(p_1\) 为真：任取其中两条相交直线确定一个平面，第三条直线分别与前两条直线相交且三线不过同一点，因此第三条直线上的两个交点是该平面内的两个不同点，第三条直线也在此平面内. 命题 \(p_2\) 为假：当三点共线时，过这条直线有无数个平面. 命题 \(p_3\) 为假：空间中不相交的两条直线还可能是异面直线. 命题 \(p_4\) 为真：若直线 \(m\perp\alpha\)，则 \(m\) 垂直于平面 \(\alpha\) 内的任意直线，故 \(m\perp l\). 所以 \(p_1\land p_4\)、\(\neg p_2\vee p_3\)、\(\neg p_3\vee\neg p_4\) 均为真，所有真命题的序号为 \(\circled{1}\)\(\circled{3}\)\(\circled{4}\). 故答案为 \(\circled{1}\)\(\circled{3}\)\(\circled{4}\).
\end{solution}

\section{解答题}

\begin{problem}
\(\triangle ABC\) 的内角 \(A\)，\(B\)，\(C\) 的对边分别为 \(a\)，\(b\)，\(c\)，已知 \(\cos^2\left( \dfrac{\pi}{2}+A \right)+\cos A=\dfrac{5}{4}\).

\begin{enumerate}
\item 求 \(A\)；
\item 若 \(b-c=\dfrac{\sqrt3}{3}a\)，证明：\(\triangle ABC\) 是直角三角形.
\end{enumerate}
\end{problem}

\begin{solution}
\begin{enumerate}
\item 因为 \(\cos\left(\dfrac\pi2+A\right)=-\sin A\)，所以 \(\cos^2\left(\dfrac\pi2+A\right)=\sin^2A=1-\cos^2A\). 令 \(u=\cos A\)，则题设化为 \(1-u^2+u=\dfrac54\)，即 \(u^2-u+\dfrac14=0\)，所以 \(\left(u-\dfrac12\right)^2=0\)，得 \(\cos A=\dfrac12\). 由于 \(0<A<\pi\)，故 \(A=\dfrac\pi3\).
\item 由正弦定理，设三角形外接圆半径为 \(R\)，则 \(a=2R\sin A=2R\sin\dfrac\pi3=\sqrt3R\). 又 \(b-c=2R(\sin B-\sin C)\)，利用和差化积及 \(B+C=\dfrac{2\pi}3\)，得 \(b-c=4R\cos\dfrac{B+C}{2}\sin\dfrac{B-C}{2}=2R\sin\dfrac{B-C}{2}\). 由题设 \(b-c=\dfrac{\sqrt3}{3}a=R\)，所以 \(\sin\dfrac{B-C}{2}=\dfrac12\). 因为 \(B,C>0\) 且 \(B+C=\dfrac{2\pi}3\)，有 \(-\dfrac\pi3<\dfrac{B-C}{2}<\dfrac\pi3\)，故 \(\dfrac{B-C}{2}=\dfrac\pi6\). 从而 \(B-C=\dfrac\pi3\)，联立 \(B+C=\dfrac{2\pi}3\)，得 \(B=\dfrac\pi2\)，\(C=\dfrac\pi6\). 因此 \(\triangle ABC\) 是直角三角形.
\end{enumerate}
\end{solution}

\begin{problem}
某沙漠地区经过治理，生态系统得到很大改善，野生动物数量有所增加. 为调查该地区某种野生动物的数量，将其分成面积相近的 \(200\) 个地块，从这些地块中用简单随机抽样的方法抽取 \(20\) 个作为样区，调查得到样本数据 \((x_i,y_i)\)（\(i=1\)，\(2\)，\(\cdots\)，\(20\)），其中 \(x_i\) 和 \(y_i\) 分别表示第 \(i\) 个样区的植物覆盖面积（单位：公顷）和这种野生动物的数量，并计算得 \(\sum_{i=1}^{20}x_i=60\)，\(\sum_{i=1}^{20}y_i=1200\)，\(\sum_{i=1}^{20}(x_i-\bar x)^2=80\)，\(\sum_{i=1}^{20}(y_i-\bar y)^2=9000\)，\(\sum_{i=1}^{20}(x_i-\bar x)(y_i-\bar y)=800\).

\begin{enumerate}
\item 求该地区这种野生动物数量的估计值（这种野生动物数量的估计值等于样区这种野生动物数量的平均数乘以地块数）；
\item 求样本 \((x_i,y_i)\)（\(i=1\)，\(2\)，\(\cdots\)，\(20\)）的相关系数（精确到 \(0.01\)）；
\item 根据现有统计资料，各地块间植物覆盖面积差异很大. 为提高样本的代表性以获得该地区这种野生动物数量更准确的估计，请给出一种你认为更合理的抽样方法，并说明理由.
\end{enumerate}

附：相关系数
\(r=\dfrac{\sum_{i=1}^{n}(x_i-\bar x)(y_i-\bar y)}{\sqrt{\sum_{i=1}^{n}(x_i-\bar x)^2\sum_{i=1}^{n}(y_i-\bar y)^2}}\)，
\(\sqrt2\approx 1.414\).
\end{problem}

\begin{solution}
\begin{enumerate}
\item 样区中动物数量的平均数为 \(\bar y=\dfrac1{20}\sum_{i=1}^{20}y_i=\dfrac{1200}{20}=60\). 按题意，该地区动物数量的估计值为样区平均数乘以地块数，因此估计值为 \(60\times200=12000\).
\item 代入相关系数公式，得 \(r=\dfrac{800}{\sqrt{80\times9000}}\). 因为 \(80\times9000=720000=(600\sqrt2)^2\)，所以 \(r=\dfrac{800}{600\sqrt2}=\dfrac{2\sqrt2}{3}\approx\dfrac{2\times1.414}{3}\approx0.94\). 因而相关系数约为 \(0.94\).
\item 可采用分层随机抽样：先按各地块的植物覆盖面积将 \(200\) 个地块分成若干层，再在各层内进行随机抽样，并按各层地块数占总体地块数的比例分配样本数. 由于各地块间植物覆盖面积差异很大，分层后每层内部的差异相对较小，且各类地块都能被抽到，样本的代表性更强，从而能够减小抽样误差，使动物数量的估计更加准确.
\end{enumerate}
\end{solution}

\begin{problem}
已知椭圆 \(C_1:\dfrac{x^2}{a^2}+\dfrac{y^2}{b^2}=1\)（\(a>b>0\)）的右焦点 \(F\) 与抛物线 \(C_2\) 的焦点重合，\(C_1\) 的中心与 \(C_2\) 的顶点重合. 过 \(F\) 且与 \(x\) 轴垂直的直线交 \(C_1\) 于 \(A\)，\(B\) 两点，交 \(C_2\) 于 \(C\)，\(D\) 两点，且 \(\lvert CD\rvert=\dfrac{4}{3}\lvert AB\rvert\).

\begin{enumerate}
\item 求 \(C_1\) 的离心率；
\item 若 \(C_1\) 的四个顶点到 \(C_2\) 的准线距离之和为 \(12\)，求 \(C_1\) 与 \(C_2\) 的标准方程.
\end{enumerate}
\end{problem}

\begin{solution}
\begin{enumerate}
\item 设椭圆的右焦点为 \(F(c,0)\)，其中 \(c=\sqrt{a^2-b^2}>0\). 因为抛物线的焦点为 \(F\)，顶点为椭圆中心 \((0,0)\)，所以抛物线开口向右，方程可设为 \(y^2=4cx\). 直线 \(x=c\) 与椭圆相交时，代入椭圆方程得 \(\dfrac{c^2}{a^2}+\dfrac{y^2}{b^2}=1\). 由 \(c^2=a^2-b^2\)，得 \(y^2=\dfrac{b^4}{a^2}\)，所以 \(\lvert AB\rvert=\dfrac{2b^2}{a}\). 直线 \(x=c\) 与抛物线相交时，得 \(y^2=4c^2\)，所以 \(\lvert CD\rvert=4c\). 由题设 \(4c=\dfrac43\cdot\dfrac{2b^2}{a}\)，从而 \(\dfrac ca=\dfrac{2b^2}{3a^2}\). 令离心率 \(e=\dfrac ca\)，又 \(b^2=a^2-c^2=a^2(1-e^2)\)，于是 \(e=\dfrac23(1-e^2)\)，即 \(2e^2+3e-2=0\). 因为 \(0<e<1\)，故 \(e=\dfrac12\).
\item 由 \(e=\dfrac12\)，得 \(c=\dfrac a2\). 抛物线的准线为 \(x=-c\). 椭圆的四个顶点为 \((a,0)\)，\((-a,0)\)，\((0,b)\)，\((0,-b)\)，它们到准线的距离之和为 \((a+c)+(a-c)+c+c=2a+2c=3a\). 由题设 \(3a=12\)，得 \(a=4\)，从而 \(c=2\)，\(b^2=a^2-c^2=16-4=12\)，即 \(b=2\sqrt3\). 所以 \(C_1\) 与 \(C_2\) 的标准方程分别为 \(\dfrac{x^2}{16}+\dfrac{y^2}{12}=1\) 和 \(y^2=8x\).
\end{enumerate}
\end{solution}

\begin{problem}
如图，已知三棱柱 \(ABC-A_1B_1C_1\) 的底面是正三角形，侧面 \(BB_1C_1C\) 是矩形，\(M\)，\(N\) 分别为 \(BC\)，\(B_1C_1\) 的中点，\(P\) 为 \(AM\) 上一点. 过 \(B_1C_1\) 和 \(P\) 的平面交 \(AB\) 于 \(E\)，交 \(AC\) 于 \(F\).

\begin{enumerate}
\item 证明：\(AA_1\parallel MN\)，且平面 \(A_1AMN\perp\) 平面 \(EB_1C_1F\)；
\item 设 \(O\) 为 \(\triangle A_1B_1C_1\) 的中心. 若 \(AO=AB=6\)，\(AO\parallel\) 平面 \(EB_1C_1F\)，且 \(\angle MPN=\dfrac{\pi}{3}\)，求四棱锥 \(B-EB_1C_1F\) 的体积.
\end{enumerate}

\bitmapfigure[float=inline,max width=6.0cm]{img/2020/national_paper_2_liberal/q20_fig1.png}
\end{problem}

\begin{solution}
\begin{enumerate}
\item 设 \(BC\) 的中点为 \(M\)，则 \(MN\) 是上下底面对应中点的连线. 由棱柱的性质，\(AA_1\parallel BB_1\parallel CC_1\)，又 \(MN\parallel BB_1\)，所以 \(AA_1\parallel MN\). 在等边三角形 \(ABC\) 中，\(AM\perp BC\)；又因为侧面 \(BB_1C_1C\) 是矩形，故 \(BB_1\perp BC\)，从而 \(AA_1\perp BC\). 因此平面 \(A_1AMN\) 内相交的两条直线 \(AM\)，\(AA_1\) 都垂直于 \(BC\)，所以 \(BC\perp\text{平面 }A_1AMN\). 又 \(B_1C_1\parallel BC\)，且 \(B_1C_1\subset\text{平面 }EB_1C_1F\)，故平面 \(A_1AMN\perp\text{平面 }EB_1C_1F\).
\item 建立空间直角坐标系，使 \(M\) 为原点，\(BC\) 为 \(x\) 轴，底面 \(ABC\) 为 \(xOy\) 平面，并设正三角形边长为 \(s\). 于是可取
\(B\left(-\dfrac s2,0,0\right)\)，\(C\left(\dfrac s2,0,0\right)\)，\(A\left(0,\dfrac{\sqrt3s}{2},0\right)\). 设棱柱的平移向量为 \(\overrightarrow{AA_1}=\overrightarrow{BB_1}=\overrightarrow{CC_1}=(0,u,h)\)，其中 \(h>0\). 设 \(P=(0,p,0)\)，则
\(A_1=(0,\dfrac{\sqrt3s}{2}+u,h)\)，\(B_1=(-\dfrac s2,u,h)\)，\(C_1=(\dfrac s2,u,h)\)，\(N=(0,u,h)\). 由于上底面也是正三角形，其中心为 \(O=(0,u+\dfrac{\sqrt3s}{6},h)\). 因而 \(\overrightarrow{AO}=(0,u-\dfrac{\sqrt3s}{3},h)\).

平面 \(EB_1C_1F\) 由直线 \(B_1C_1\) 和点 \(P\) 确定. 取该平面内的两个方向向量 \((1,0,0)\) 和 \(P-B_1=(\dfrac s2,p-u,-h)\)，可得其一个法向量为 \((0,h,p-u)\). 由 \(AO\parallel\text{平面 }EB_1C_1F\)，有 \((0,u-\dfrac{\sqrt3s}{3},h)\cdot(0,h,p-u)=0\)，即 \(h(p-\dfrac{\sqrt3s}{3})=0\). 所以 \(p=\dfrac{\sqrt3s}{3}\).

又由 \(AO=AB=6\)，得 \(s=6\)，且 \(\lvert AO\rvert=\lvert PN\rvert=6\). 因此 \(p=2\sqrt3\)，并且 \(\overrightarrow{PM}=(0,-2\sqrt3,0)\)，\(\overrightarrow{PN}=(0,u-2\sqrt3,h)\). 由 \(\angle MPN=\dfrac\pi3\)，得
\(\dfrac{\overrightarrow{PM}\cdot\overrightarrow{PN}}{\lvert PM\rvert\lvert PN\rvert}=\dfrac12\)，即 \(\dfrac{-2\sqrt3(u-2\sqrt3)}{2\sqrt3\times6}=\dfrac12\). 解得 \(u-2\sqrt3=-3\). 再由 \(\lvert PN\rvert=6\)，得 \(h^2=36-9=27\)，所以 \(h=3\sqrt3\).

此时平面 \(EB_1C_1F\) 的方程为 \(\sqrt3y+z-6=0\). 它与底面 \(z=0\) 的交线为 \(y=2\sqrt3\)，故 \(EF\parallel BC\). 在底面内，\(AM=3\sqrt3\)，而 \(AP=\sqrt3\)，所以由相似三角形 \(\triangle AEF\sim\triangle ABC\)，得 \(EF=\dfrac{AP}{AM}BC=2\). 又 \(B_1C_1=BC=6\)，且 \(PN\) 垂直于这两条平行边，\(PN=6\)，所以四边形 \(EB_1C_1F\) 的面积为 \(\dfrac{2+6}{2}\times6=24\). 点 \(B\) 到平面 \(EB_1C_1F\) 的距离为 \(\dfrac{\lvert\sqrt3\cdot0+0-6\rvert}{\sqrt{(\sqrt3)^2+1^2}}=3\). 因此四棱锥的体积为 \(\dfrac13\times24\times3=24\).
\end{enumerate}
\end{solution}

\begin{problem}
已知函数 \(f(x)=2\ln x+1\).

\begin{enumerate}
\item 若 \(f(x)\le 2x+c\)，求 \(c\) 的取值范围；
\item 设 \(a>0\)，讨论函数 \(g(x)=\dfrac{f(x)-f(a)}{x-a}\) 的单调性.
\end{enumerate}
\end{problem}

\begin{solution}
\begin{enumerate}
\item 要使不等式对定义域内一切 \(x>0\) 成立，需有 \(2\ln x+1\leq2x+c\)，即 \(c\geq2\ln x-2x+1\). 令 \(h(x)=2\ln x-2x+1\)，则 \(h'(x)=\dfrac2x-2\). 当 \(0<x<1\) 时 \(h'(x)>0\)，当 \(x>1\) 时 \(h'(x)<0\)，所以 \(h(x)\) 在 \(x=1\) 处取得最大值 \(h(1)=-1\). 因而 \(c\geq-1\) 时不等式恒成立，且由 \(x=1\) 可知 \(c\) 不可能小于 \(-1\)，所以 \(c\in[-1,+\infty)\).
\item 定义域为 \(x>0\) 且 \(x\ne a\). 化简得 \(g(x)=\dfrac{2\ln x-2\ln a}{x-a}=\dfrac{2\ln(x/a)}{x-a}\). 求导得 \(g'(x)=\dfrac{2\left(1-\dfrac ax-\ln\dfrac xa\right)}{(x-a)^2}\). 令 \(t=\dfrac xa>0\)，分子中的函数为 \(\varphi(t)=1-\dfrac1t-\ln t\)，且 \(\varphi'(t)=\dfrac{1-t}{t^2}\). 因此 \(\varphi\) 在 \((0,1)\) 上递增、在 \((1,+\infty)\) 上递减，并且 \(\varphi(1)=0\)，从而 \(\varphi(t)<0\)（\(t\ne1\)）. 所以 \(g'(x)<0\)，故 \(g(x)\) 在 \((0,a)\) 和 \((a,+\infty)\) 上均单调递减.
\end{enumerate}
\end{solution}

\section{选考题：解答题}

\begin{problem}
已知曲线 \(C_1\) 的参数方程为
\(\begin{cases}
x=4\cos^2\theta,\\
y=4\sin^2\theta
\end{cases}\)（\(\theta\) 为参数），\(C_2\) 的参数方程为
\(\begin{cases}
x=t+\dfrac{1}{t},\\
y=t-\dfrac{1}{t}
\end{cases}\)（\(t\) 为参数）.

\begin{enumerate}
\item 将 \(C_1\)，\(C_2\) 的参数方程化为普通方程；
\item 以坐标原点为极点，\(x\) 轴正半轴为极轴建立极坐标系. 设 \(C_1\)，\(C_2\) 的交点为 \(P\)，求圆心在极轴上，且经过极点和 \(P\) 的圆的极坐标方程.
\end{enumerate}
\end{problem}

\begin{solution}
\begin{enumerate}
\item 对曲线 \(C_1\)，由 \(x=4\cos^2\theta\)，\(y=4\sin^2\theta\)，得 \(x+y=4\)，且 \(x\geq0\)，\(y\geq0\). 因此 \(C_1\) 的普通方程为 \(x+y=4\)（\(x\geq0\)，\(y\geq0\)）. 对曲线 \(C_2\)，参数 \(t\ne0\)，由参数方程消去 \(t\) 得 \(x^2-y^2=\left(t+\dfrac1t\right)^2-\left(t-\dfrac1t\right)^2=4\)，所以 \(C_2\) 的普通方程为 \(x^2-y^2=4\).
\item 联立两曲线的普通方程，并考虑 \(C_1\) 的范围，得 \(y=4-x\)，以及 \(x^2-(4-x)^2=4\). 化简得 \(8x-16=4\)，所以 \(x=\dfrac{5}{2}\)，\(y=\dfrac{3}{2}\)，即 \(P\left(\dfrac{5}{2},\dfrac{3}{2}\right)\). 设所求圆的圆心为 \((h,0)\)，因为圆经过极点 \(O(0,0)\)，半径为 \(\lvert h\rvert\)，其直角坐标方程为 \((x-h)^2+y^2=h^2\)，即 \(x^2+y^2=2hx\). 将点 \(P\) 代入，得 \(\dfrac{25}{4}+\dfrac94=2h\cdot\dfrac52\)，从而 \(h=\dfrac{17}{10}\). 令极坐标为 \((\rho,\theta)\)，代入 \(x=\rho\cos\theta\)，\(x^2+y^2=\rho^2\)，得 \(\rho^2=\dfrac{17}{5}\rho\cos\theta\)，所以所求圆的极坐标方程为 \(\rho=\dfrac{17}{5}\cos\theta\).
\end{enumerate}
\end{solution}

\begin{problem}
已知函数 \(f(x)=\lvert x-a^2\rvert+\lvert x-2a+1\rvert\).

\begin{enumerate}
\item 当 \(a=2\) 时，求不等式 \(f(x)\ge 4\) 的解集；
\item 若 \(f(x)\ge 4\)，求 \(a\) 的取值范围.
\end{enumerate}
\end{problem}

\begin{solution}
\begin{enumerate}
\item 当 \(a=2\) 时，\(f(x)=\lvert x-4\rvert+\lvert x-3\rvert\). 分三种情况讨论：当 \(x\leq3\) 时，\(f(x)=7-2x\)，由 \(7-2x\geq4\) 得 \(x\leq\dfrac32\)；当 \(3\leq x\leq4\) 时，\(f(x)=1<4\)，无解；当 \(x\geq4\) 时，\(f(x)=2x-7\)，由 \(2x-7\geq4\) 得 \(x\geq\dfrac{11}{2}\). 所以不等式的解集为 \(\left(-\infty,\dfrac32\right]\cup\left[\dfrac{11}{2},+\infty\right)\).
\item 令 \(u=a^2\)，\(v=2a-1\). 对任意实数 \(x\)，由三角不等式得 \(f(x)=\lvert x-u\rvert+\lvert x-v\rvert\geq\lvert u-v\rvert=\lvert a^2-2a+1\rvert=(a-1)^2\). 当 \(x\) 位于 \(u\) 与 \(v\) 之间时取等号，所以 \(f(x)\geq4\) 恒成立的充要条件是 \((a-1)^2\geq4\). 解得 \(a\leq-1\) 或 \(a\geq3\)，故 \(a\) 的取值范围为 \(\left(-\infty,-1\right]\cup[3,+\infty)\).
\end{enumerate}
\end{solution}
